%%=============================================================================
%% Implementatie en configuratie van NetBox Discovery
%%=============================================================================

\chapter{\IfLanguageName{dutch}{Implementatie en configuratie van NetBox Discovery}{Implementation and configuration of NetBox Discovery}}%
\label{ch:fase4}

\section{Configuratie van de Orb agent voor Diode}

In dit gedeelte wordt de concrete configuratie van de Orb agent beschreven die gebruikt wordt om discovery-data via Diode naar NetBox te sturen. De belangrijkste configuratie gebeurt in YAML-bestanden, aangevuld met shellscripts die de vereiste omgevingsvariabelen instellen en de agentcontainer starten.

\subsection{Voorbeeldconfiguratie voor device discovery}

\begin{listing}[H]
  \caption{Voorbeeldconfiguratie van de Orb agent voor device discovery via Diode}
  \label{lst:orb-device-discovery-config}
  \begin{minted}[fontsize=\small]{yaml}
orb:
  # Zet het globale logniveau van de agent omhoog
  log_level: DEBUG 
  config_manager:
    active: local

  backends:
    common:
      diode:
        target: grpc://localhost:8080/diode
        client_id: ${DIODE_CLIENT_ID}
        client_secret: ${DIODE_CLIENT_SECRET}
        agent_name: discovery_agent_01
        dry_run: true
        dry_run_output_dir: /opt/orb
    device_discovery:
      host: 0.0.0.0
      port: 8072
      # Dit vertelt de backend wrapper om op debug te loggen
      log_level: DEBUG
      log_format: TEXT

  policies:
    device_discovery:
      oostkamp_router:
        config:
          
          schedule: "*/30 * * * *" 
          options:
            port_scan_ports: [8722]
            port_scan_timeout: 1.0
            
            debug: true 
          defaults:
            site: "Oostkamp Ridefort"
            role: "Router"
            tenant: "Oostkamp"
        scope:
          - hostname: "10.23.4.225"
            driver: "ros"
            username: "${DEVICE_USERNAME}"
            password: "${DEVICE_PASSWORD}"
            optional_args:
              port: 8722
              use_ssl: false
              debug: true 
  \end{minted}
\end{listing}

\subsection{Voorbeeldconfiguratie voor SNMP discovery}

\begin{listing}[H]
  \caption{Voorbeeldconfiguratie van de Orb agent voor SNMP discovery}
  \label{lst:orb-snmp-discovery-config}
  \begin{minted}[fontsize=\small]{yaml}
orb:
  config_manager:
    active: local

  backends:
    common:
      diode:
        target: grpc://127.0.0.1:8080/diode
        client_id: ${DIODE_CLIENT_ID}
        client_secret: ${DIODE_CLIENT_SECRET}
        agent_name: discovery_agent_01

  policies:
    snmp_discovery:
      oostkamp_ridefort:
        config:
          schedule: "*/30 * * * *"
          timeout: 300
          snmp_timeout: 5
          retries: 3
          lookup_extensions_dir: "/opt/orb/snmp-extensions/lookup_extensions"
          defaults:
            site: "Oostkamp Ridefort"
            role: "network"
            rack: "Oostkamp Ridefort"
            tags: ["snmp-discovery", "orb"]
            device:
              tenant: "Oostkamp"
              rack: "oostkamp_ridefort"
              description: "SNMP discovered device"

        scope:
          targets:
            - host: "10.23.4.225"
              override_defaults:
                role: "Router"

            - host: "10.23.4.226"
              override_defaults:
                role: "Switch"

            - host: "10.23.4.231-236"
              override_defaults:
                role: "Access Point"

          authentication:
            protocol_version: "SNMPv2c"
            community: "public"
  \end{minted}
\end{listing}

\subsection{Shellscripts voor genereren van configuratie en starten van de agent}

De volgende shellscripts genereren op basis van OAuth2-credentials de definitieve \texttt{agent.yaml}-configuratie en starten de Orb agentcontainer.

\begin{listing}[H]
  \caption{Shellscript voor het genereren van \texttt{agent.yaml} voor SNMP discovery}
  \label{lst:orb-snmp-script}
  \begin{minted}[fontsize=\small]{bash}
#!/usr/bin/env bash
set -e

SCRIPT_DIR="$(cd "$(dirname "${BASH_SOURCE[0]}")" && pwd)"
cd "$SCRIPT_DIR"

mkdir -p "$SCRIPT_DIR/snmp-extensions"

CREDS_FILE="/opt/diode/oauth2/client/client-credentials.json"
ALT_CREDS_FILE="/vagrant/diode/oauth2/client/client-credentials.json"

if [ ! -f "$CREDS_FILE" ]; then
  if [ -f "$ALT_CREDS_FILE" ]; then
    CREDS_FILE="$ALT_CREDS_FILE"
  else
    echo "Fout: $CREDS_FILE niet gevonden."
    exit 1
  fi
fi

if ! command -v envsubst >/dev/null 2>&1; then
  echo "envsubst niet gevonden. Installeer met: sudo apt-get install -y gettext-base"
  exit 1
fi

export DIODE_CLIENT_ID="diode-ingest"
export DIODE_CLIENT_SECRET=$(sudo jq -r '.[] | select(.client_id == "diode-ingest") | .client_secret' "$CREDS_FILE")

if [ -z "$DIODE_CLIENT_SECRET" ]; then
  echo "Fout: kon diode-ingest secret niet ophalen."
  exit 1
fi

export DEVICE_USERNAME="user"
export DEVICE_PASSWORD="pass"

envsubst '${DIODE_CLIENT_ID} ${DIODE_CLIENT_SECRET} ${DEVICE_USERNAME} ${DEVICE_PASSWORD}' < agent.yaml.snmp > agent.yaml
echo "agent.yaml gegenereerd met credentials."

exec sudo docker run --rm --network host \
  -v "$SCRIPT_DIR:/opt/orb" \
  -e PIP_BREAK_SYSTEM_PACKAGES=1 \
  netboxlabs/orb-agent:latest run -c /opt/orb/agent.yaml
  \end{minted}
\end{listing}

\begin{listing}[H]
  \caption{Shellscript voor het genereren van \texttt{agent.yaml} voor router discovery}
  \label{lst:orb-router-script}
  \begin{minted}[fontsize=\small]{bash}
#!/usr/bin/env bash
set -e

SCRIPT_DIR="$(cd "$(dirname "${BASH_SOURCE[0]}")" && pwd)"
cd "$SCRIPT_DIR"

mkdir -p "$SCRIPT_DIR/snmp-extensions"

CREDS_FILE="/opt/diode/oauth2/client/client-credentials.json"
ALT_CREDS_FILE="/vagrant/diode/oauth2/client/client-credentials.json"

if [ ! -f "$CREDS_FILE" ]; then
  if [ -f "$ALT_CREDS_FILE" ]; then
    CREDS_FILE="$ALT_CREDS_FILE"
  else
    echo "Fout: $CREDS_FILE niet gevonden."
    exit 1
  fi
fi

if ! command -v envsubst >/dev/null 2>&1; then
  echo "envsubst niet gevonden. Installeer met: sudo apt-get install -y gettext-base"
  exit 1
fi

export DIODE_CLIENT_ID="diode-ingest"
export DIODE_CLIENT_SECRET=$(sudo jq -r '.[] | select(.client_id == "diode-ingest") | .client_secret' "$CREDS_FILE")

if [ -z "$DIODE_CLIENT_SECRET" ]; then
  echo "Fout: kon diode-ingest secret niet ophalen."
  exit 1
fi

export DEVICE_USERNAME="user"
export DEVICE_PASSWORD="pass"

envsubst '${DIODE_CLIENT_ID} ${DIODE_CLIENT_SECRET} ${DEVICE_USERNAME} ${DEVICE_PASSWORD}' < agent.yaml.router > agent.yaml
echo "agent.yaml gegenereerd met credentials."

exec sudo docker run --rm --network host \
  -v "$SCRIPT_DIR:/opt/orb" \
  -e INSTALL_DRIVERS_PATH=/opt/orb/drivers-router.txt \
  -e PIP_BREAK_SYSTEM_PACKAGES=1 \
  netboxlabs/orb-agent:2.5.0 run -c /opt/orb/agent.yaml
  \end{minted}
\end{listing}

\section{Werking in de praktijk en aanvullende automatisering}

\subsection{Wat de Orb-agent wel oplevert}

De SNMP discovery en de device discovery via de Orb agent functioneren naar verwachting: apparaten worden ontdekt en via Diode in NetBox geregistreerd. De configuratie in de voorgaande secties is in de praktijk ingezet en levert een betrouwbare basis voor inventarisatie.

\subsection{Beperking: LLDP-buren in NetBox}

LLDP-buren (\emph{neighbors}) worden via deze keten niet automatisch als relaties in NetBox aangemaakt. Ondanks een correcte device discovery (met onder meer capabilities voor topologie en buren in de configuratie) bleek het niet mogelijk om die LLDP-informatie consistent automatisch te laten landen in NetBox. Dit vormt een beperking voor volledig geautomatiseerde kabel- en buurrelaties op basis van alleen de standaard NetBox Discovery-workflow.

\subsection{Aanvulling: Streamlit-toepassing en NetBox API voor kabel- en LLDP-data}

LLDP-buren (\emph{neighbors}) komen niet automatisch als relaties in NetBox terecht via de Orb-agent-discoveryflow. Om toch kabel- en topologierelaties uit te werken, is een eigen webtoepassing gebouwd met Streamlit (Python). De kern van deze aanpak is:

\begin{itemize}
  \item een apparaat naar keuze uitlezen (per vendor/protocol) om o.a. interfaces, IP-informatie en LLDP-neighbors te verzamelen;
  \item de verzamelde data via de NetBox REST API schrijven naar de relevante NetBox objecten (device, interfaces, IP-adressen en kabels).
\end{itemize}

Voor de apparaten waarvoor een geschikte driver beschikbaar is, levert deze pipeline een praktisch alternatief: waar NetBox Discovery/Orb de LLDP-relaties niet automatisch aanmaakt, haalt Streamlit de LLDP-informatie op en maakt het vervolgens kabelobjecten op basis van de LLDP-neighbors.

\subsubsection{Agent Manager tab: Orb-agent starten vanuit Streamlit}

De eerste tab van de applicatie beheert de Orb agent via Docker. Op basis van de gekozen modus (`SNMP`, `Router` of `Switch`) selecteert Streamlit automatisch het juiste YAML-configuratiebestand en startscript.

\begin{listing}[H]
  \caption{Streamlit: selectie van Orb-agentconfig en startscript}
  \label{lst:streamlit-agent-manager}
  \begin{minted}[fontsize=\small]{python}
AGENT_MAP = {
    "SNMP": {"config": "agent.yaml.snmp", "script": "start-agent-snmp.sh"},
    "Router": {"config": "agent.yaml.router", "script": "start-agent-router.sh"},
    "Switch": {"config": "agent.yaml.switch", "script": "start-agent-switch.sh"}
}

mode = st.radio("Selecteer Agent:", ["SNMP", "Router", "Switch"], horizontal=True)
cfg_file = AGENT_MAP[mode]["config"]
run_script = AGENT_MAP[mode]["script"]

if st.button(f"Start Agent ({mode})"):
    log_file = os.path.join(AGENT_DIR, "agent_output.log")
    with open(log_file, "a") as f:
        subprocess.Popen(
            ["bash", run_script],
            cwd=AGENT_DIR,
            stdout=f,
            stderr=f
        )
    st.success(f"Start-commando voor {mode} verzonden.")
    st.rerun()
  \end{minted}
\end{listing}

\subsubsection{Network Scanner tab: bulk scan en NetBox API import}

De tweede tab scant toestellen (of een IP-range) en genereert daarna optioneel een reeks CSV-exports (achtergrond). Vervolgens kan je vanuit dezelfde interface bulkdata pushen naar NetBox via de NetBox API.

\begin{listing}[H]
  \caption{Streamlit: scan-form en bulk API push}
  \label{lst:streamlit-api-import}
  \begin{minted}[fontsize=\small]{python}
with st.form("scan_form"):
    vendor = st.selectbox("Merk", ["MikroTik", "Huawei"])
    ip_input = st.text_input("IP-adres of Range").strip()
    site = st.text_input("NetBox Site", value="Oostkamp").strip()

    user = st.text_input("Gebruikersnaam", value="sa_netbox_orb").strip()
    password = st.text_input("Wachtwoord", type="password").strip()
    role = st.selectbox("NetBox Role", ["Router", "Switch"])
    rack = st.text_input("NetBox Rack", value="").strip()
    tenant = st.text_input("NetBox Tenant", value="").strip()

    submitted = st.form_submit_button("Start Scan & Genereer CSV's")

if submitted:
    target_ips = parse_ip_input(ip_input)
    # per IP: ophalen van feiten + genereren van CSV-exports (afhankelijk van vendor)
    ...

with st.form("api_push_form"):
    api_options = st.multiselect(
        "Selecteer modules om te pushen:",
        ["Device", "Interfaces", "IP's", "Cables"],
        default=["Device", "Interfaces", "IP's", "Cables"]
    )
    nb_url = st.text_input("NetBox URL", value="---").strip()
    nb_token = st.text_input("NetBox API Token", type="password", value="---").strip()
    push_btn = st.form_submit_button(
        f"Bevestig: Push {len(st.session_state.scan_results)} apparaten naar NetBox"
    )

    if push_btn:
        for d in st.session_state.scan_results:
            success, message = sync_device_to_netbox(
                ip=d["ip"], username=d["user"], password=d["password"],
                site_name=d["site"], role_name=d["role"], rack_name=d["rack"],
                tenant_name=d["tenant"], custom_name=d["naam"],
                nb_url=nb_url, nb_token=nb_token,
                options=api_options, status_box=status,
                vendor=d["vendor"],
            )
  \end{minted}
\end{listing}

\subsubsection{LLDP-neighbors ophalen (MikroTik) en kabels aanmaken in NetBox}

In onderstaande (verkorte) codefragmenten wordt zichtbaar hoe je Streamlit de LLDP-neighbors verkrijgt en daarna kabels creëert via de NetBox API.

\begin{listing}[H]
  \caption{Extract: LLDP-neighbors ophalen en voorbereid verwerken in \texttt{sync\_device\_to\_netbox}}
  \label{lst:streamlit-sync-lldp}
  \begin{minted}[fontsize=\small]{python}
def sync_device_to_netbox(ip, username, password, site_name, role_name, rack_name, tenant_name, custom_name, nb_url, nb_token, options, status_box, vendor):
    try:
        status_box.write(f"⏳ Verbinden met {vendor} ({ip}) via NAPALM...")
        
        # --- DATA OPHALEN (LOKAAL OF VIA BRUG) ---
        if vendor == "MikroTik":
            driver = get_network_driver('ros')
            device = driver(ip, username, password, optional_args={'port': 8722})
            device.open()
            
            facts = device.get_facts()
            status_box.write("📡 Data ophalen van de router...")
            interfaces = device.get_interfaces() if ("Interfaces" in options or "IP's" in options or "Cables" in options) else {}
            ips = device.get_interfaces_ip() if "IP's" in options else {}
            
            lldp_neighbors = {}
            if "Cables" in options:
                status_box.write("🔌 LLDP (bekabeling) informatie ophalen...")
                try:
                    lldp_neighbors = device.get_lldp_neighbors()
                except Exception:
                    pass 
            device.close()
            
    except Exception as e:
        return False, f"Fout tijdens NetBox API sync: {str(e)}"
  \end{minted}
\end{listing}

\begin{listing}[H]
  \caption{Extract: kabels aanmaken in NetBox op basis van LLDP-neighbors}
  \label{lst:streamlit-create-cables}
  \begin{minted}[fontsize=\small]{python}
        # D. CABLES
        if "Cables" in options and nb_dev:
            count_cables = 0
            for local_iface, neighbors in lldp_neighbors.items():

                if "eoip" in local_iface.lower() or "bridge" in local_iface.lower():
                    continue

                for neighbor in neighbors:
                    remote_host = neighbor.get('hostname', '').strip()
                    remote_port = str(neighbor.get('port', '')).strip()

                    if not remote_host:
                        continue

                    if "eoip" in remote_port.lower() or "bridge" in remote_port.lower():
                        continue

                    if re.match(r'^([0-9a-fA-F]{4}-){2}[0-9a-fA-F]{4}$', remote_port):
                        remote_port = "eth1"

                    intf_a_search = list(nb.dcim.interfaces.filter(device_id=nb_dev.id, name=local_iface))
                    if not intf_a_search:
                        intf_a_search = list(nb.dcim.interfaces.filter(device_id=nb_dev.id, name=expand_port_name(local_iface)))

                    if not intf_a_search:
                        continue

                    nb_intf_a = intf_a_search[0]

                    if getattr(nb_intf_a, 'cable', None):
                        continue

                    remote_devs = list(nb.dcim.devices.filter(name=remote_host))
                    if len(remote_devs) == 0 or len(remote_devs) > 1:
                        continue
                    nb_dev_b = remote_devs[0]

                    remote_intfs = list(nb.dcim.interfaces.filter(device_id=nb_dev_b.id, name=remote_port))
                    if len(remote_intfs) == 0:
                        remote_intfs = list(nb.dcim.interfaces.filter(device_id=nb_dev_b.id, name=expand_port_name(remote_port)))

                    if len(remote_intfs) == 0 or len(remote_intfs) > 1:
                        continue
                    nb_intf_b = remote_intfs[0]

                    try:
                        nb.dcim.cables.create(
                            a_terminations=[{"object_type": "dcim.interface", "object_id": nb_intf_a.id}],
                            b_terminations=[{"object_type": "dcim.interface", "object_id": nb_intf_b.id}],
                            status="connected"
                        )
                        count_cables += 1
                    except Exception:
                        pass

            messages.append(f"✅ {count_cables} nieuwe kabels verbonden via LLDP-detectie.")
  \end{minted}
\end{listing}

\subsection{Ruckus access points versus overige apparaten}

Voor Ruckus access points bestaat geen NAPALM-driver. Deze toestellen worden daarom eerst in NetBox gebracht via SNMP discovery (zoals beschreven in de SNMP-configuratie hierboven). Voor die AP's blijft de aanvulling via de Streamlit-pipeline beperkt tot wat via SNMP of handmatige stappen mogelijk is; de nadruk van de custom scripts ligt op andere platformen.

\subsection{Huawei-switch en MikroTik-router}

Na de initiële registratie in NetBox worden met de Streamlit-toepassing gericht de Huawei-switch en de MikroTik-router verder uitgelezen. Afhankelijk van de vendor worden de relevante gegevens opgehaald en via de NetBox API ingevoerd (o.a. IP-adressen en \emph{cable connections}). Zo ontstaat een samenhangend beeld van topologie en bekabeling voor de kerncomponenten, ter aanvulling op wat SNMP discovery en device discovery automatisch leveren.

